% Автоматически перенесено из КОНТЕКСТ_ВЕРСИЯ_3.docx

\chapter*{Каталонские графства в XI--XIII~вв.}
\addcontentsline{toc}{chapter}{Каталонские графства в XI--XIII~вв.}

\section*{§1. Формирование политической автономии графства Барселонского}
\addcontentsline{toc}{section}{§1. Формирование политической автономии графства Барселонского}

В 801~г. в ходе длительной осады войска франкского короля Людовика Благочестивого, сына Карла Великого, овладели Барселоной, ранее находившейся под властью мусульманского эмира\footnote{\emph{Арский И.В. }Очерки по истории средневековой Каталонии до соединения с Арагоном. (VIII--XII века). Л.: Ленингр. гос. ун-т, 1941. С.~6.}. Город был включён в состав Испанской марки (пограничной каролингской области, выполнявшей функцию буферной зоны между Франкской державой и аль-Андалусом) и стал административным центром региона.

Заключение в 857~г. мирного соглашения между вали Сарагосы и графом Барселоны закрепило христиано-мусульманскую границу по р.~Льобрегат, что на несколько столетий предопределило различия в политическом и социальном развитии северных и южных областей Каталонии\footnote{Название Каталония появилось относительно поздно и датируется приблизительно 2-й половиной XI-- началом XII~вв.: \emph{Арский~И.В.} Указ. соч. С.~6; \emph{Мильская}\emph{~Л.Т}. Очерки из истории деревни в Каталонии X-XII~вв. М.: Изд-во Акад. наук СССР, 1962. С.~4; \emph{Варьяш}\emph{~}\emph{И.И., Астахов М.А.} Арагон, Каталония и Наварра в VIII--XI~вв. \emph{// }История Испании / под ред. \emph{В.А.~Ведюшкина}, \emph{Г.А.~Поповой}. М.: «Индрик», 2012. Т.~1. С.~219--220; \emph{Sabat}\emph{é }\emph{i}\emph{ }\emph{Curull}\emph{~}\emph{F}\emph{.} L’origen medieval de la identitat catalana // Anàlisi històrica de la identitat catalana. [Barcelona]: Institut d’Estudis Catalans, 2015. P.~25; \emph{Salrach J.M.} Orígens i transformacions de la senyoria a Catalunya (segles IX--XIII) // Revista d’historia medieval. 1997. №~8. P.~28 и др.}. Постепенное ослабление Каролингов в X~в. способствовало укреплению положения местной династии\footnote{\emph{Варьяш}\emph{ И.И., Астахов М.А.} Арагон, Каталония и Наварра в VIII--XI~вв. \emph{// }История Испании / под ред. \emph{В.А.~Ведюшкина}, \emph{Г.А.~Поповой}. М.: «Индрик», 2012. Т.~1. С.~221.}. Сочетание пограничного положения, политической раздробленности и династической консолидации сформировало условия для раннего правового плюрализма и устойчивости локальных норм\footnote{\emph{Bisson T.N.} Feudalism in Twelfth-Century Catalonia // Publications de l’École française de Rome. 1980. Т.~44. №~1. P.~186.}.

В последующий период значительную роль в укреплении графской власти сыграл маркграф \textbf{Винифред}\textbf{ Мохнатый} \textbf{(878--898)}\footnote{Здесь и далее имена графов Барселонских и правителей Барселонского дома передаются преимущественно в каталанской форме, а не в латинизированной или кастильской традиции: Рамон Беренгер (а не Раймунд Беренгарий), Жауме (а не Хайме), Пере (а не Педро или Пётр) и т.д.}, которого традиционно считали основателем Барселонской династии\footnote{\emph{Варьяш}\emph{ И.И., Астахов М.А.} Арагон, Каталония и Наварра в VIII--XI~вв. \emph{// }История Испании / под ред. \emph{В.А.~Ведюшкина}, \emph{Г.А.~Поповой}. М.: «Индрик», 2012. Т.~1. С.~219--220.}. Закрепление наследственного характера власти графов Барселонских и усиление её связей с местной аристократией (баронами) создали предпосылки для постепенного ослабления зависимости от каролингского контроля.

Несмотря на включение в состав государства Каролингов, Барселона на протяжении IX--X~вв. сохраняла статус пограничного христианского форпоста. Это проявилось особенно отчётливо в 985~г., когда город и окрестные христианские монастыри подверглись разгрому войсками аль-Мансура. Отсутствие эффективной помощи со стороны франкского короля стало важным фактором, позволившим \textbf{Борелю II}\textbf{ }\textbf{(954--992) }отказаться от признания каролингского верховенства и закрепить фактическую самостоятельность графства Барселонского\footnote{Там же. С.~222.}.

К XII~в. земли будущего каталонского Принципата (кат. \emph{Principat}\emph{ de }\emph{Catalunya})\footnote{\emph{Л.Т.~Мильская} отмечала, что этимология названия Каталония весьма спорна, и появляется это обозначение не ранее XI~в.: \emph{Мильская}\emph{~}\emph{Л.Т.} К вопросу о структуре господствующего класса в Каталонии X--XII~вв. // Средние века. Вып.~47. М.: Наука, 1984. С.~22.} представляли собой совокупность графств (Барселона, Осона, Жирона, Уржель, Серданья-Конфлент-Бесалу, Пальярс, Рибагорса, Ампурьяс, Руссильон и др.), сформировавшихся на основе каролингской пограничной организации, но уже обладавших собственной политико-правовой идентичностью.

С возвышением графа Барселонского \textbf{Рамона }\textbf{Беренгера}\textbf{ I Старого (1035--1076)} над другими каталонскими графами за Барселонской династией закрепилось старшинство. Вокруг неё постепенно сложилась система сеньориальных связей, объединявшая прочие графства региона и обеспечивавшая координацию их политических и военных действий. Рамон Беренгер I и его третья супруга Альмодис де ла Марш (1052--1071) также проводили активную внешнюю политику\footnote{\emph{Варьяш}\emph{~И.И., Астахов~М.А.} Каталония и Арагон в XI--XIV~вв. \emph{// }История Испании / под ред. \emph{В.А.~Ведюшкина}, \emph{Г.А.~Поповой}. М.: «Индрик», 2012. Т.~1. С.~261; \emph{Aurell M.} Les noces du comte: mariage et pouvoir en Catalogne (785--1213). Paris: Publications de la Sorbonne, 1995. P.~271.}. В ходе конфликтов и переговоров с мусульманскими правителями северо-востока Пиренейского полуострова граф добился установления регулярных даннических платежей от некоторых тайф\footnote{\textbf{Тайфы} (араб. \textbf{\arabictext{طائفة}} (\emph{{\unicodefallback ṭ}ā{\unicodefallback ʾ}ifa}) --- «группа», «сторона», «партия»)--- мусульманские государственные образования, возникшие на территории аль-Андалуса в результате распада Кордовского халифата и со временем достигшие фактической самостоятельности. См. \emph{Варьяш}\emph{~И.И.} Очерк политической истории \emph{// }История Испании / под ред. \emph{В.А.~Ведюшкина}, \emph{Г.А.~Поповой}. М.: «Индрик», 2012. Т.~1. С.~135; Taifa // Gran Enciclopèdia Catalana. Enciclopèdia.cat. URL: \url{https://www.enciclopedia.cat/gran-enciclopedia-catalana/taifa-0} (дата обращения: 16.05.2026).}. Параллельно Барселонская династия укрепила свои позиции к северу от Пиренеев: связи с Разесом и Каркассоном демонстрировали включённость каталонской знати в политическое пространство Окситании и расширяли сеть династических и вассальных взаимодействий\footnote{\emph{Варьяш}\emph{ И.И., Астахов М.А.} Арагон, Каталония и Наварра в VIII--XI~вв. \emph{// }История Испании / под ред. \emph{В.А.~Ведюшкина}, \emph{Г.А.~Поповой}. М.: «Индрик», 2012. Т.~1. С.~222;}.

\section*{§2. Территориальная экспансия XII~в. и интеграция новых земель}
\addcontentsline{toc}{section}{§2. Территориальная экспансия XII~в. и интеграция новых земель}

В первой половине XII~в. Бесалу, Серданья и часть Ампурьяса окончательно вошли в состав владений барселонских графов. Прованс оказался под их властью лишь на определённое время. Одновременно укрепилась автономия каталонской церкви, чему способствовало восстановление архиепископской кафедры в Таррагоне (1118) и ослабление зависимости от Нарбоннской архиепископии\footnote{Там же. С.~265; \emph{Филиппов И.С.} Средиземноморская Франция в раннее средневековье: Проблема становления феодализма. М.: Скрипторий, 2000. С.~201.}.

Ключевой этап политической трансформации пришёлся на правление \textbf{Рамона }\textbf{Беренгера}\textbf{ IV (1131--1162)}. Заключённое в 1137~г. соглашение о браке с наследницей арагонского престола Петронилой Арагонской (1136--1173), было окончательно оформлено их свадьбой в 1150~г. Она положила начало династической унии графства Барселоны с королевством Арагон и сформировало монархию, в рамках которой правитель объединял различные титулы и юрисдикции при сохранении правовых традиций отдельных территорий\footnote{\emph{Bisson T.N.} Feudalism in Twelfth-Century Catalonia // Publications de l’École française de Rome. 1980. Т.~44. №~1. P.~186; \emph{Idem.} The Problem of Feudal Monarchy: Aragon, Catalonia, and France // Speculum. 1978. Т.~53, Vol.~3. P.~464.}.

Правление Рамона Беренгера IV ознаменовалось решающим продвижением христиан к югу: отвоевание Тортосы (1148) и Льейды (1149) завершило включение в сферу каталонского господства большей части территорий современной Каталонии (за исключением Валенсии и Балеарских островов)\footnote{\emph{Idem.} The Problem of Feudal Monarchy: Aragon, Catalonia, and France // Speculum. [Medieval Academy of America, Cambridge: Cambridge University Press, Chicago: University of Chicago Press], 1978. Т.~53, Vol.~3. P.~467.}.При этом в правовом плане отвоёванные территории требовали перераспределения земли, фиксации повинностей, разграничения юрисдикций и определения статуса нехристианского населения\footnote{\emph{Guinot }\emph{i}\emph{ Rodríguez E.} The expansion of a European feudal monarchy during the 13\textsuperscript{th} Century: the Catalan-Aragonese Crown and the consequences of the conquest of the kingdoms of Majorca and Valencia // Catalan Historical Review. 2009. Núm. 2. P.~37--38.}.

Земли, расположенные к югу от р.~Льобрегат, в историографии принято называть Новой Каталонией. Их освоение сопровождалось переселением населения из северных графств, что уже к концу XII~в. обеспечило устойчивое каталонское присутствие в регионе\footnote{\emph{Pacheco Catalán N.} Ut sis statoret habitator: La colonización feudal de la Catalunya Nova~: la inmigración occitana (siglos XII--XIII). Granada: Universidad de Granada, 2025. P.~22.}. Противопоставление Старой и Новой Каталонии восходит к традиции XIII~в. и впервые зафиксировано у \textbf{Пере Альберта}\footnote{\textbf{Пере Альберт (1233}\textbf{ -- }\textbf{ок}\textbf{. }\textbf{1261)} --- каталонский юрист, предположительно получивший образование в Болонском университете, каноник Собора Святого Креста и Святой Евлалии. Известен как автор трактата по праву «Коммеморации», в котором он попытался соединить локальные записи обычного права с римским правовым наследием. Подробнее о нём и его труде см.: \emph{Ferran i Planas E.} El jurista Pere Albert i les «Commemoracions» / ed. T.~de Montagut i~Estragués. Barcelona: Institut d’Estudis Catalans. 2006. P.~172.}, который относил к Новой Каталонии территории к югу от р.~Льобрегат, включённые в состав владений барселонских графов в ходе завоеваний середины XII~в.\footnote{Цит. по: \emph{Ferran i Planas E.} El jurista Pere Albert i les «Commemoracions» / ed. T. de Montagut~i~Estragués. Barcelona: Institut d’Estudis Catalans, 2006. P.~172.}. В современной историографии (Ф.~Сабате-и-Куруй и А.~Виргили-и-Колет) это деление интерпретируется как отражение различий в хронологии и характере феодализации каталонских территорий\footnote{\emph{Sabaté i Curull F.} La noció d’Espanya a la Catalunya medieval // Acta historica et archaeologica mediaevalia. 1998. P.~375--390; \emph{Virgili i Colet A.} Gent nova. La colonització feudal de la Catalunya Nova (segles XII--XIII) // Butlletí de la Societat Catalana d’Estudis Històrics. 2010. Т.~21. P.~100.}.

\section*{§3. Королевская власть и сословно-представительные институты в Каталонии XIII~в.}
\addcontentsline{toc}{section}{§3. Королевская власть и сословно-представительные институты в Каталонии XIII~в.}

После кризиса начала XIII~в., когда после поражения \textbf{Пере II Католика }\textbf{(1174--1213)} в сражении при Мюре (1213), власть перешла к его малолетнему сыну Жауме~I, при котором произошла стабилизация политического порядка и развитие сословного представительства\footnote{\emph{Варьяш}\emph{~И.И., Астахов~М.А.} Каталония и Арагон в XI--XIV~вв. \emph{// }История Испании / под ред. \emph{В.А.~Ведюшкина}, \emph{Г.А.~Поповой}. М.: «Индрик», 2012. Т.~1. С.~277.}. Уже в первой половине XIII~в. каталонские графства функционировали как целостное политико-правовое пространство, объединённое под властью Барселонского графского дома\footnote{Там же. С.~278--279.}.

Внутренняя политика \textbf{Жауме I Завоевателя (1213--1276) }строилась на согласовании интересов основных сословных групп --- знати, духовенства и горожан, поскольку каталонская монархия не обладала централизованным аппаратом управления, сопоставимым с кастильским\footnote{Там же. С.~278--279.}. Король выступал прежде всего как арбитр и гарант традиционных норм.

В этих условиях возрастало значение сословно-представительных институтов, прежде всего кортесов Каталонии, где согласовывались налоговые обязательства и военные повинности и подтверждались права различных земель и общин. Уже к середине XIII~в. сложилась практика, при которой налогообложение требовало согласия сословий, что существенно ограничивало возможность односторонних решений короля\footnote{Там же. С.~281.}.

Каталонская знать XIII~в. сохраняла значительную автономию и собственную судебную власть, однако её положение постепенно менялось. В отличие от XI--XII~вв., когда крупные сеньоры выступали почти независимыми центрами политической силы, в XIII~в. они всё более интегрировались в структуру Арагонской Короны через участие в кортесах, военных кампаниях и управлении новыми территориями\footnote{Там же. С.~280.}. Стоит также отметить, что именно в правление Жауме I Завоевателя действие «Барселонских обычаев» распространилось на все каталонские графства.

Правовая культура знати сочетала традиционную феодальную автономию с признанием юрисдикции Короны. Наиболее заметные изменения внутренней структуры Каталонии XIII~в. были связаны с ростом городов и их политической роли. Барселона, Таррагона, Льейда и другие центры активно развивались благодаря торговле и колонизации, что усиливало экономический потенциал региона и позиции городских общин. Городам и даже небольшим селениям, пребывавшим на момент пожалования привилегий в состоянии хозяйственного запустения, предоставлялись и подтверждались привилегии, закреплявшие их самоуправление, судебные полномочия местной администрации и торговые права\footnote{\emph{Фрязинов С.В.} Материалы для изучения фуэрос и поселенных хартий // Средние века. Вып.~22. М.: Наука, 1962. С.~163.} Через муниципальные советы и консулов городские общины всё активнее включались в политическую жизнь короны, а их представительство в кортесах способствовало оформлению горожан как особой сословной группы\textsuperscript{ }\footnote{\emph{Кучумов В.А. }Кортесы Арагона и Каталонии: генезис и специфика // Проблемы испанской истории. М.: Наука, 1992. C.~146.}. Внутреннее право городов формировалось как совокупность статутов, обычаев и королевских привилегий, что усиливало плюралистический характер правовой культуры Каталонии.

Городская экономика требовала фиксации договоров, имущественных прав и корпоративных норм, что ускоряло внедрение письменных форм правового регулирования и элементов общеевропейской правовой традиции. Участие каталонской знати и горожан в освоении завоёванных территорий усиливало их позиции внутри Короны и расширяло социальную базу монархии\footnote{\emph{Варьяш}\emph{~И.И., Астахов~М.А.} Каталония и Арагон в XI--XIV~вв. \emph{// }История Испании / под ред. \emph{В.А.~Ведюшкина}, \emph{Г.А.~Поповой}. М.: «Индрик», 2012. Т.~1. С.~280.}.

Показательный пример --- кортесы 1283~г. в Барселоне, созванные Пере~III~Арагонским в условиях острой потребности короны в политической и финансовой поддержке\footnote{\emph{Кучумов}\emph{~}\emph{В.А. }Кортесы Арагона и Каталонии: генезис и специфика // Проблемы испанской истории. М.: Наука, 1992. C.~146.}. На них духовенство, знать и представители городских общин выступили как участники переговоров с монархом: они представили петиции с требованиями прав. Именно по итогам работы этих кортесов, Барселоне была пожалована привилегия «Знатные мужи Барселоны признали…» (Recognoverunt proceres)\footnote{Подробнее см. раздел «Источниковедение» настоящей работы.}.

\section*{§4. Система управления}
\addcontentsline{toc}{section}{§4. Система управления}

Управление каталонскими землями в XII--XIII~вв. формировалась в условиях расширения территории и усложнения политической структуры графств. Как и другие королевства Пиренейского полуострова королевская власть в Каталонии осуществлялась посредством сети представителей правителя и локальных должностных лиц, действовавших в рамках сеньориальных отношений и множественных юрисдикций\footnote{\emph{Варьяш}\emph{~И.И., Астахов~М.А.} Каталония и Арагон в XI--XIV~вв. \emph{// }История Испании / под ред. \emph{В.А.~Ведюшкина}, \emph{Г.А.~Поповой}. М.: «Индрик», 2012. Т.~1. С.~282. Например, часть Тортосы после её отвоевания находилась в юрисдикции духовно-рыцарских орденов (тамплиеров, и наиболее важным должностным лицом для них являлся магистр ордена). \emph{Варьяш}\emph{~}\emph{И.И.} Потестарные институты и должности в Испании в V--XV~вв. // Властные институты и должности в Европе в Средние века и раннее Новое время. М.: КДУ, 2011. С.~417.}.

Ключевыми фигурами этой структуры выступали \textbf{вегеры} (лат. \emph{vicarius}, кат. \emph{veguer}), представлявшие верховную власть графа или короля в территориальных округах --- вегериях. Вегеры председательствовали в местных судах, обеспечивали поддержание общественного порядка, организовывали оборону и сбор налогов, а также следили за исполнением королевских распоряжений\footnote{\emph{Варьяш}\emph{ И.И.} Указ. соч. С.~428--429.}. \textbf{Батлы}\textbf{ }(лат. baiulus, кат.~\emph{batlle}) помогали вегерам контролировать налогообложение и исполнение повинностей на местах и могли рассматривать имущественные или хозяйственные споры, а также были наделены судебными полномочиями, кроме того, как отмечает Е.~Родон-Бинуэ, батлы могли управлять владениями отдельных феодалов\footnote{Там же. С.~433. Cм. также \emph{Rodon Binue E.} El lenguaje técnico del feudalismo en el siglo XI en Cataluña. (Coniribucrón al estudio del latin medieval). Barcelona: Escuela de filología, 1957. P.~37. Батлов, как нам кажется, некорректно определять как отдельную социальную прослойку между знатью и простолюдинами, как это было сделано \emph{Л.Т.~Мильской}: \emph{Мильская}\emph{~Л.Т.} К вопросу о структуре господствующего класса в Каталонии X--XII~вв. // Средние века. Вып.~47. М.: Наука, 1984. С.~25. По нашему мнению, этот термин указывает на административную функцию, а не социальную принадлежность конкретного лица.}. \textbf{С}\textbf{айоны}\textbf{ }(кат.~\emph{saig}\emph{,} \emph{saiones}\emph{, }\emph{sayones}) обеспечивали вызов сторон в суд, исполнение решений суда, а также участвовали в сборе штрафов и иных платежей\footnote{\emph{Варьяш}\emph{~И.И.} Потестарные институты и должности в Испании в V--XV~вв. // Властные институты и должности в Европе в Средние века и раннее Новое время. М.: КДУ, 2011. С.~418.}.

Отправление правосудия в каталонских землях обеспечивалось также широким кругом специализированных должностных лиц --- \textbf{судей} (\emph{iudices}), \textbf{писцов} и \textbf{нотариев}. Они работали при дворе правителя, при вегерах и в городских судах, применяя местные обычаи, привилегии и нормы писанного права\footnote{\emph{Sabat}\emph{é }\emph{i}\emph{ }\emph{Curull}\emph{ }\emph{F}\emph{.} ~El Veguer a Catalunya. Anàlisi del funcionament de la jurisdicció reial al segle XIV // Butlletí de la Societat Catalana d’Estudis Històrics. 1995. Т.~6. P.~149;\emph{ }\emph{Варьяш}\emph{~}\emph{И.И.} Указ. соч. С.~418.}.

В обязанности \textbf{кастелянов} (лат. \emph{castlani}) входили оборона территории, надзор за населением и исполнение судебных решений. Замок оставался базовой единицей власти и юрисдикции, вокруг которой строилась организация сельского управления\footnote{\emph{Мильская}\emph{~Л.Т.} К вопросу о структуре господствующего класса в Каталонии X--XII~вв. // Средние века. Вып.~47. М.: Наука, 1984. С.~25; \emph{Bonnassie}\emph{ Р. }Ор. cit. T.~2. P.~785.}.

Существенную роль в административно-правовом устройстве Каталонии играли и коллективные институты --- собрания знати и баронов, а также церковные советы и синоды. Именно на них утверждались положения «мира Божьего», ограничивавшие частные войны, а также принимались решения о фиксации и подтверждении местных обычаев и привилегий\footnote{\emph{Bisson}~\emph{T.}\emph{N}\emph{.} The Organized Peace in Southern France and Catalonia (c. 1140--1223) // The American Historical Review. 1977. Vol.~82. P.~291. Подробнее о существующих в историографии концепциях «мира и перемирия Божьего» см.: \emph{Gonzalvo }\emph{i}\emph{ Bou}\emph{~}\emph{G}. Les constitucions de Pau i Treva de Catalunya: (segles XI -- XIII). Barcelona, 1994. P.~XIX--XXXII.}. Участие знати и духовенства в подобных ассамблеях придавало нормам коллективную легитимацию.

Церковные структуры, прежде всего епископские курии и капитулы, также играли административную роль, участвуя в судебных процедурах, арбитраже и хранении документов. Церковь обладала собственной юрисдикцией, но одновременно взаимодействовала со светской властью, формируя общее правовое пространство региона\footnote{\emph{Варьяш}\emph{~И.И., Астахов~М.А.} Каталония и Арагон в XI--XIV~вв. \emph{// }История Испании / под ред. \emph{В.А.~Ведюшкина}, \emph{Г.А.~Поповой}. М.: «Индрик», 2012. Т.~1. С.~269.}.

Городские советы и консулы осуществляли судебные и управленческие функции, взаимодействуя с вегерами и батлами как с представителями королевской власти. Наглядным примером этого является городская администрация Барселоны: в XIII~в.: она находилась под управлением вегера, которому помогал батл, также действовал \textbf{Г}\textbf{ородской совет,} в состав которого входила преимущественно знать (кат. \emph{Consell}\emph{ de }\emph{ple})\footnote{Там же. С.~281.}. В 1258~г. Жауме~I~Завоеватель учредил посты \textbf{4}\textbf{ }\textbf{вегеров}\textbf{ с 8 советниками}. Кроме того, городской совет теперь включал около двухсот человек, а в 1265~г. был преобразован в \textbf{Совет ста} (кат.~\emph{Consell}\emph{ de }\emph{Cent})\footnote{\emph{Font }\emph{i Rius}\emph{~}\emph{J.M.} Estudis sobre els drets i institucions locals en la Catalunya medieval. Barcelona: Edicions Universitat Barcelona, 1985. P.~475.}. Этот орган обладал правом чеканки монеты и назначал представителей Барселоны за её пределами. Также в Барселоне действовал \textbf{Морско}\textbf{й}\textbf{ совет} (кат.~\emph{Consolat}\emph{ de }\emph{Mar}), объединявший корабельщиков и купцов, деятельность которых была тесно связана с морской торговлей\footnote{\emph{Ferrer Mallol}\emph{~}\emph{M.T. }El Consolat de Mar i els consolats d’Ultramar, instrument i manifestació de l’expansió del comerç català // L'Expansió Catalana a la Mediterrània a la Baixa Edat Mitjana: actes del séminaire-seminari, celebrat a Barcelona el 20 d'abril de 1998 / coord. por María Teresa Ferrer Mallol, Damien Coulon,. Barcelona: Consejo Superior de Investigaciones Científicas (CSIC), 1999. P.~59.}.

\section*{§5. Социальные группы и их правовое положение}
\addcontentsline{toc}{section}{§5. Социальные группы и их правовое положение}

Социальная структура Каталонии XII--XIII~вв. определяется тем, что статус того или иного лица определялся его местом в иерархии земельных держаний. Родовая знать (магнаты и бароны), в том числе потомки графских родов, владели крупными феодами и замками и участвовали в собраниях графской курии\footnote{\emph{Мильская}\emph{~}\emph{Л}\emph{.}\emph{Т}\emph{.} К вопросу о структуре господствующего класса в Каталонии X--XII~вв. // Средние века. Вып.~47. М.: Наука, 1984. С.~25; \emph{Bonnassie }\emph{Р}\emph{. }Ор. cit. T.~2. P.~785.}. Верхушку каталонской знати составляли \textbf{местные графы}, \textbf{виконты} и \textbf{кастеляны} наиболее важных замков, названные в «Барселонских обычаях» как «магнатами земли» (лат.~\emph{magnates}\emph{ }\emph{terrae}\emph{)}\footnote{\emph{Мильская}\emph{~}\emph{Л.Т.} Указ. соч. С.~24.}. Ниже располагались \textbf{комиторы}\textbf{ }(лат. \emph{comitores}),\textbf{ }потомки побочных линий высшей знати, кастелянов и вегеров, и \textbf{вальвассоры}\textbf{ }(лат. \emph{vasvassores}), занимавшие промежуточное положение между высшей знатью и рыцарством\footnote{\emph{Альтамира}\emph{-и-}\emph{Кревеа}\emph{~}\emph{Р.} История средневековой Испании. СПб.: Евразия, 2003. Т.~1. С.~282.}. \textbf{Рыцари} (лат. \emph{milites}) составляли нижний и наиболее многочисленный уровень среди привилегированных категорий населения и были напрямую связаны с несением военной службы своим сеньорам. Как показал И.С.~Филиппов, в Каталонии, по крайней мере с XII~в., \textbf{бакалари}\textbf{ями} (лат. \emph{baccalarii}) называли \emph{«крестьянских парней, взятых в свиту сеньора»}\footnote{\emph{Филиппов~И.С.} Средиземноморская Франция в раннее средневековье: Проблема становления феодализма. М.: Скрипторий, 2000. С.~497.}, которые сопровождали господина в военных походах и сражались рядом с ним, но не могли достичь положения рыцаря.

Городское население в источниках обозначалось двумя терминами (лат.~\emph{cives}, лат. \emph{burgenses}\emph{)}, которые указывали на принадлежность человека к городской общине, но не проводили строгого различия между разными группами горожан.

Наиболее широким обозначением сельского населения был термин \textbf{крестьяне} (лат. \emph{rustici}), не сводившийся только к лично зависимым людям, известным как \textbf{ременсы} (кат. \emph{remences}). В Старой Каталонии, где феодализация завершилась раньше всего, к XIII~в. сложился комплекс крестьянских повинностей, позднее получивший название \textbf{«дурны}\textbf{х}\textbf{ обыча}\textbf{ев}\textbf{»} (кат.~\emph{mals}\emph{ }\emph{usos})\footnote{\emph{Пискорский}\emph{~}\emph{В.К.} Вопрос о значении и происхождении шести «дурных обычаев» в Каталонии: С прил. неизд. документов. Киев: тип. Ун-та св. \emph{Владимира Н.Т.}~Корчак-Новицкого, 1899. С.~8--9.}. Эти обычаи закрепляли личную зависимость крестьянина от сеньора и рассматривали его держание и имущество как включённые в сферу власти сеньора при сохранении традиционных платежей и повинностей.

К числу классических «дурных обычаев» относили:

\textbf{Ременса} (кат. \emph{remen}\emph{ç}\emph{a}\emph{ }\emph{personal}) --- выкуп, который крестьянин обязан был выплатить при переходе с земли или смене сеньора; именно от этой повинности происходило само обозначение зависимых крестьян;

\textbf{Интестия} (лат. \emph{intestia}) --- обычай, предусматривавший передачу сеньору лучшей головы скота или части имущества крестьянина, умершего без завещания;

\textbf{Экзоркия} (лат. \emph{exorquia}) --- право сеньора на выморочное имущество крестьянина, умершего без наследников;

\textbf{Кугусия} (лат. \emph{cugutia}) --- право сеньора на долю имущества крестьянина в случае супружеской измены его жены, трактуемой как оскорбление чести сеньора;

\textbf{Арсия} (лат. arcia) --- штраф в пользу сеньора при пожаре в крестьянском хозяйстве, независимо от вины держателя;

\textbf{Фирма }\textbf{д’эсполи}\textbf{ }\textbf{форсада} (лат. ferma de spoli)--- обязательный платёж господину при заключении брака крестьянина или передаче хозяйства\footnote{\emph{Варьяш}\emph{~}\emph{И.И., Астахов}\emph{~}\emph{М.А.} Арагонская Корона в XV~в. \emph{// }История Испании / под ред. \emph{В.А.~Ведюшкина}, \emph{Г.А.~Поповой}. М.: «Индрик», 2012. Т.~1. С.~346. П.~Боннасси настаивал на том, что включение собственно «ременсы» в состав «дурных обычаев» некорректно: \emph{Bonnassie}\emph{~}\emph{Р. }Ор. cit. T.~2. P.~824.}.

Иногда к «дурным обычаям» относят и так называемое «право первой ночи», однако источники классического Средневековья практически не подтверждают его реального применения: в позднесредневековый период оно существовало лишь в символической форме, связанной с брачными церемониями и подтверждением власти сеньора\footnote{\emph{Варьяш}\emph{ И.И., Астахов М.А.} Арагонская Корона в XV~в. \emph{// }История Испании / под ред. \emph{В.А.~Ведюшкина}, \emph{Г.А.~Поповой}. М.: «Индрик», 2012. Т.~1. С.~346.}. В Старой Каталонии «дурные обычаи» были широко распространены, поскольку там феодальные отношения оформились раньше, чем на отвоёванных у мавров землях\footnote{\emph{Bisson T.N.} Feudalism in Twelfth-Century Catalonia // Publications de l’École française de Rome. 1980. Т.~44. Vol~1. P.~186.}. В Новой Каталонии зависимость крестьян имела менее жестокий характер в связи с необходимостью хозяйственного освоения земель и стремлением власти привлечь поселенцев на отвоёванные территории\footnote{\emph{Мильская}\emph{ Л.Т.} Указ. соч. С.~27.}.

Социальную структуру Каталонии дополняли многочисленные \textbf{общины иноверцев}, мусульман и иудеев\footnote{В «Барселонских обычаях» также встречается термин «еретики»: Usatges de Barcelona. Us. 60 (64): «...christiani et sarraceni, iudei et heretici...». P.~96. Под «еретиками» в «Барселонских обычаях» едва ли следует понимать одну строго определённую группу. В каталонском обществе XII--XIII~вв. это обозначение относилось к катарам, или альбигойцам. Довольно подробная библиография по этому вопросу: \emph{Thouzellier C. }Catharisme et valdéisme en Languedoc: à la fin du XIIe et au début du XIIIe siècle: politique pontificale, controverses. Paris: Presses universitaires de France, 1965. P.~14.}. По подсчётам М.Т.~Феррер-Майоль, в Тортосе и Льейде в XIII~в. 25--33\% жителей являлись мусульманами\footnote{\emph{Ferrer Mallol M.T.} L'Aljama islámica de Tortosa a la Baixa Edat Mitjana // Recerca. Centre d'Estudis Històrics del Baix Ebre, 2003. Núm. 7. P.~186.}. Иудеи, общины которых были сосредоточены в Барселоне, Жироне, Таррагоне и Льейде, составляли до 10\% городского населения, однако играли весьма важную в торговле, финансах и управлении\footnote{\emph{Curto Homedes A}. Topografia del call jueu de Tortosa // Recerca. 1999. Núm. 3. P.~10.}.

И мусульмане, и иудеи сохраняли внутреннюю правовую автономию и собственные судебные институты, подчиняясь при этом фискальной и политической власти Арагонской Короны. При составлении грамот и разрешении споров иноверцы нередко обращались к переводчикам и посредникам. Таким образом, их язык, вера и правовое положение отличались от христианского, однако составление документа подчинялось нормам, установленным христианской властью\footnote{\emph{Варьяш}\emph{ И.И., Зеленина Г.С., Попова Г.А. }Три веры \emph{// }История Испании / под ред. \emph{В.А.~Ведюшкина}, \emph{Г.А.~Поповой}. М.: «Индрик», 2012. Т.~1. С.~188.}.
